\section{Visualisierung des Erosionsprozesses}

Die nachfolgenden vier Abbildungen stellen den theoretischen Wasserabsaug- und Erosionsprozess dar, der zur Entstehung des Grand Canyon führte. Jede Visualisierung beleuchtet unterschiedliche Aspekte des Modells: von zweidimensionalen Querschnitten über dreidimensionale Perspektiven bis hin zur mathematischen Formalisierung und der zeitlichen Prozessabfolge.

\subsection{Zweidimensionale Querschnittsdarstellung (erosion\_2d.svg)}

Die erste Visualisierung präsentiert neun aufeinanderfolgende Querschnitte, die den stufenweisen Erosionsprozess über einen Zeitraum von etwa 220 Tagen dokumentieren. Jeder der neun Sub-Plots repräsentiert einen diskreten Zyklus im Gesamtprozess und vermittelt ein detailliertes Bild der progressiven Canyon-Formation.

\subsubsection*{Sub-Plot 1: Zyklus 0 (Tag 0--5) -- Ausgangszustand}

Der erste Querschnitt zeigt den initialen Zustand unmittelbar nach Ende der Sedimentationsphase. Die horizontale Achse erstreckt sich über eine Breite von 10 Kilometern (von $-5$ km bis $+5$ km), während die vertikale Achse Tiefen bis zu 2.160 Metern darstellt. 

Der Untergrund wird durch eine einheitliche braune Fläche (Farbe: \#8B7355) dargestellt, die das frisch abgelagerte, noch nicht verfestigte Sedimentgestein symbolisiert. Über diesem Untergrund liegt eine durchgehende Wasserschicht in charakteristischem Blau (Farbe: \#4A90E2, Transparenz 70\%), die die gesamte Oberfläche bedeckt und bis zum Nullniveau reicht.

Eine Textannotation im oberen Bereich des Plots hebt hervor: \textit{„Ausgangszustand -- Wasser vollständig über dem Boden"}. Diese Phase entspricht dem biblischen Tag 150, an dem das Gewässer sein Maximum erreichte. Es findet noch keine Erosion statt; das System ist in einem metastabilen Gleichgewichtszustand. Die Legende in der oberen rechten Ecke identifiziert drei Komponenten: Gestein (braun), erodierter Bereich (heller Sand, \#D4A574) und Wasser (blau).

\subsubsection*{Sub-Plot 2: Zyklus 1 (Tag 5--10) -- Erster Erosionszyklus}

Der zweite Querschnitt markiert den Beginn des katastrophalen Rücklaufprozesses. Deutlich erkennbar ist die erste Absenkung des Wasserspiegels: Die blaue Wasserfläche hat sich um etwa 200 Meter nach unten verschoben. 

Im Zentrum des Querschnitts (bei $x = 0$) beginnt sich eine V-förmige Erosionsstruktur auszubilden. Die Erosionstiefe beträgt in dieser Phase bereits etwa 225 Meter. Der erodierte Bereich wird in einem helleren Sandton (Farbe: \#D4A574, Transparenz 60\%) dargestellt, der sich vom dunkleren, unerodiertem Gestein abhebt.

Eine gestrichelte horizontale Linie bei $z = 0$ markiert das ursprüngliche Bodenniveau, während eine zweite gestrichelte Linie bei etwa $-225$ m die erste Erosionsstufe kennzeichnet. Diese Linien sind in Schwarz mit 30\% Transparenz gehalten und haben eine Linienbreite von 0,8 Punkten.

Das Verhältnis von Breite zu Tiefe der Erosionsrinne beträgt in dieser Phase etwa 1:0,008, was der im Hauptdokument hergeleiteten Geometrie entspricht. Die Textannotation \textit{„Erosionstiefe: 225 m"} erscheint zentral im erodierten Bereich in einem halbtransparenten weißen Rahmen.

\subsubsection*{Sub-Plot 3: Zyklus 2 (Tag 10--15) -- Zweite Erosionsphase}

Der dritte Querschnitt zeigt eine deutliche Vertiefung des Erosionsprozesses. Der Wasserstand ist weiter auf etwa 1.400 Meter über dem ursprünglichen Bodenniveau abgesunken, was einem Rückgang von etwa 44\% entspricht.

Die V-förmige Canyon-Struktur hat sich nun auf eine Tiefe von etwa 450 Metern erweitert. Zwei horizontale gestrichelte Linien markieren die beiden bisherigen Erosionsstufen bei $z = -225$ m und $z = -450$ m. Diese Linien visualisieren die charakteristischen „Rillen" des Grand Canyon, die durch die stufenweise Erosion entstehen.

Der erodierte Bereich zeigt nun eine deutliche Schichtung: Die oberste Schicht (erste Erosionsphase) erscheint geringfügig dunkler als die zweite, was auf unterschiedliche Verfestigungsgrade hindeutet. Die Erosionsbreite an der Oberfläche beträgt nun etwa 3,6 Kilometer, während die maximale Tiefe im Zentrum bei 450 Metern liegt.

Die Wassermasse über dem erodierten Bereich wird in leuchtendem Blau dargestellt und zeigt durch ihre reduzierte Höhe den fortschreitenden Rücklaufprozess. Der Hintergrund des Plots ist in einem warmen Beige (Farbe: \#E8DCC8) gehalten, was den Wüstencharakter der späteren Landschaft andeutet.

\subsubsection*{Sub-Plot 4: Zyklus 3 (Tag 15--20) -- Mittlere Erosionsphase}

Im vierten Querschnitt manifestiert sich die kontinuierliche Vertiefung des Canyons. Die Erosion hat nun eine Tiefe von 675 Metern erreicht, was etwa 37,5\% der finalen Grand-Canyon-Tiefe von 1.800 Metern entspricht.

Drei gestrichelte horizontale Linien bei $-225$ m, $-450$ m und $-675$ m markieren die drei abgeschlossenen Erosionszyklen. Diese gleichmäßig vertikalen Abstände von jeweils 225 Metern illustrieren die Regelmäßigkeit des Prozesses, die durch die konstante Zyklusperiode von etwa 5 Tagen bedingt ist.

Der Wasserspiegel liegt nun bei etwa 1.125 Metern über dem ursprünglichen Boden, was einem Rückgang von 62,5\% des Anfangswasserstands entspricht. Die V-förmige Erosionsstruktur zeigt eine zunehmende Komplexität: Die Flanken des Canyons weisen leichte Unregelmäßigkeiten auf, die durch unterschiedliche Erosionsraten in verschiedenen Gesteinssorten entstehen.

Der erodierte Bereich nimmt nun einen signifikanten Teil des Querschnitts ein. Die hellsandfarbe Zone erstreckt sich von etwa $-4,5$ km bis $+4,5$ km an der Oberfläche und verengt sich konisch zur maximalen Tiefe im Zentrum. Die Textannotation gibt die präzise Erosionstiefe an: \textit{„Erosionstiefe: 675 m"}.

Das Raster (Grid) des Plots, in Grau mit 30\% Transparenz, ermöglicht eine präzise Ablesung der Dimensionen. Die Achsenbeschriftungen sind in 8-Punkt-Schrift gehalten: „Breite (km)" für die x-Achse und „Tiefe (m)" für die y-Achse.

\subsubsection*{Sub-Plot 5: Zyklus 4 (Tag 20--25) -- Fortgeschrittene Erosion}

Der fünfte Querschnitt dokumentiert den Übergang in die zweite Hälfte des Erosionsprozesses. Die Gesamterosionstiefe beträgt nun 900 Meter, genau die Hälfte der finalen Tiefe.

Vier horizontale gestrichelte Linien zeichnen die bisherigen Erosionsstufen nach. Der Wasserstand ist auf etwa 900 Meter über dem Ausgangsniveau gesunken, was bedeutet, dass die Wassersäule nun exakt der erodierten Tiefe entspricht -- ein kritischer Punkt im Modell, an dem die hydraulischen Verhältnisse sich zu ändern beginnen.

Die Erosionsbreite an der Oberfläche hat sich auf etwa 7,2 Kilometer ausgedehnt. Die Canyon-Flanken zeigen einen nahezu linearen Verlauf mit einer Neigung von etwa 1:8 (Breite zu Tiefe), was gut mit realen Grand-Canyon-Vermessungen übereinstimmt.

Der erodierte Bereich zeigt nun eine deutliche Stratifizierung mit vier erkennbaren Schichten, jede etwa 225 Meter mächtig. Diese Schichten entsprechen den vier abgeschlossenen Erosionszyklen und werden in geringfügig unterschiedlichen Helligkeitsstufen der Sandfarbe dargestellt.

Die Wassermasse über dem Canyon ist deutlich reduziert, nimmt aber immer noch einen substantiellen Teil des Plots ein. Die blaue Färbung mit 70\% Transparenz ermöglicht es, die unterlagernden Erosionsstrukturen durchscheinen zu sehen.

\subsubsection*{Sub-Plot 6: Zyklus 5 (Tag 25--30) -- Späte Erosionsphase}

Im sechsten Querschnitt erreicht die Erosion eine Tiefe von 1.125 Metern, was 62,5\% der finalen Tiefe entspricht. Fünf gestrichelte Linien markieren die fünf abgeschlossenen Erosionszyklen.

Der Wasserstand liegt nun bei etwa 675 Metern über dem ursprünglichen Boden. Bemerkenswert ist, dass der Wasserspiegel nun unterhalb der Hälfte des Ausgangsniveaus liegt, was auf eine beschleunigte Rücklaufphase hindeutet. Diese Beschleunigung ist konsistent mit dem im Hauptdokument beschriebenen quadratischen Geschwindigkeits-Druck-Zusammenhang: $v = \sqrt{2\Delta P / \rho}$.

Die V-förmige Canyon-Struktur ist nun deutlich ausgeprägt. Die Erosionsbreite an der Oberfläche beträgt etwa 9 Kilometer, während sich der Canyon zum Grund hin auf wenige hundert Meter verengt. Die Flanken zeigen eine leichte Konvexität, die durch die unterschiedlichen Erosionsraten während der verschiedenen Zyklen entsteht.

Die Textannotation \textit{„Erosionstiefe: 1.125 m"} ist zentral im Canyon platziert. Der erodierte Bereich dominiert nun das Bild, während die verbleibende Wassermasse nur noch etwa 38\% des ursprünglichen Volumens ausmacht.

\subsubsection*{Sub-Plot 7: Zyklus 6 (Tag 30--35) -- Vorletzte Phase}

Der siebte Querschnitt zeigt den vorletzten Erosionszyklus mit einer Gesamttiefe von 1.350 Metern. Sechs horizontale gestrichelte Linien markieren die Erosionsstufen in Abständen von jeweils 225 Metern.

Der Wasserstand ist auf etwa 450 Meter über dem ursprünglichen Boden gesunken, was nur noch 25\% des Ausgangswasserstands entspricht. Die Wasserschicht erscheint nun als relativ dünne blaue Linie über dem massiven Canyon.

Die Erosionsbreite an der Oberfläche nähert sich dem Maximum von etwa 10,8 Kilometern. Die Canyon-Flanken zeigen nun eine komplexe Struktur mit mehreren Terrassierungen, die den einzelnen Erosionszyklen entsprechen. Jede Terrassierung repräsentiert eine Pause zwischen den aktiven Erosionsphasen.

Der erodierte Bereich ist in sechs deutlich erkennbare Schichten stratifiziert. Die oberste Schicht (älteste Erosion) erscheint leicht dunkler als die unterste (jüngste Erosion), was den zunehmenden Verfestigungsgrad des Sediments reflektiert.

Die Textannotation gibt die Erosionstiefe präzise an. Das Raster ermöglicht eine genaue Vermessung: Die maximale Tiefe liegt bei $-1.350$ m, während sich die Oberfläche bei $z = 0$ befindet. Der Wasseroberfläche liegt bei $z = +450$ m.

\subsubsection*{Sub-Plot 8: Zyklus 7 (Tag 35--40) -- Penultimate Erosion}

Im achten Querschnitt erreicht die Erosion eine Tiefe von 1.575 Metern, was 87,5\% der finalen Tiefe entspricht. Sieben gestrichelte Linien visualisieren die sieben abgeschlossenen Erosionszyklen.

Der Wasserstand ist auf etwa 225 Meter über dem ursprünglichen Boden gesunken -- nur noch 12,5\% des Ausgangswasserstands. Die blaue Wasserschicht erscheint als schmales Band über dem Canyon und deutet das bevorstehende Ende des Rücklaufprozesses an.

Die V-förmige Canyon-Struktur ist nun nahezu vollständig ausgebildet. Die Erosionsbreite an der Oberfläche beträgt das Maximum von etwa 12,6 Kilometern. Die Canyon-Flanken zeigen eine ausgeprägte Terrassierung mit sieben deutlich erkennbaren Stufen.

Der erodierte Bereich dominiert vollständig den Querschnitt. Die Schichtung ist außerordentlich deutlich: Sieben Lagen von jeweils etwa 225 Metern Mächtigkeit sind klar voneinander abgegrenzt. Die Farbgraduierung von oben (dunkler) nach unten (heller) illustriert den zeitlichen Ablauf der Erosion.

Die Textannotation \textit{„Erosionstiefe: 1.575 m"} ist im mittleren Canyon-Bereich platziert. Das Canyon-Profil zeigt nun eine charakteristische Form, die der realen Geometrie des Grand Canyon bemerkenswert ähnlich ist.

\subsubsection*{Sub-Plot 9: Zyklus 8 (Tag 40--220) -- Endzustand}

Der neunte und letzte Querschnitt präsentiert den finalen Zustand nach Abschluss aller Erosionszyklen. Die Gesamterosionstiefe beträgt exakt 1.800 Meter -- die dokumentierte Maximalttiefe des Grand Canyon.

Das Wasser ist vollständig abgeflossen; keine blaue Schicht ist mehr vorhanden. Acht horizontale gestrichelte Linien markieren die acht Haupterosionszyklen, obwohl das reale Modell von 40--50 Zyklen ausgeht. Die Darstellung konzentriert sich auf die Hauptphasen zur besseren Visualisierung.

Die V-förmige Canyon-Struktur ist vollständig ausgebildet. Die maximale Breite an der Oberfläche beträgt etwa 14,4 Kilometer, was einem Breite-zu-Tiefe-Verhältnis von 8:1 entspricht -- konsistent mit realen Vermessungen des Grand Canyon im Bereich des South Rim.

Der erodierte Bereich zeigt eine komplexe Schichtung mit acht klar erkennbaren Lagen. Die Farbvariation von dunkelbeige (oben) zu hellsand (unten) vermittelt einen Eindruck der zeitlichen Tiefe des Prozesses. Die Canyon-Flanken weisen ausgeprägte Terrassierungen auf, die den einzelnen Erosionspausen entsprechen.

Eine hervorgehobene Textannotation in der Canyon-Mitte verkündet in Fettschrift: \textit{„Finale Tiefe: 1.800 m"}. Der gelbe Hintergrund der Textbox (Farbe: gelb, Transparenz 80\%) hebt diese Information als Höhepunkt der Sequenz hervor.

Der Titel des Sub-Plots lautet: „Zyklus 8 von 8 (Tag 40--220)", was andeutet, dass die finale Phase erheblich länger dauerte als die vorherigen Zyklen. Dies reflektiert die im Hauptdokument beschriebene Verlangsamung des Rücklaufprozesses gegen Ende, wenn der hydraulische Druck abnimmt.

Das Hintergrundfarbschema (Beige, \#E8DCC8) vermittelt den Eindruck einer trockenen, ariden Landschaft -- der moderne Grand Canyon. Das Koordinatengitter ermöglicht präzise Messungen: Die x-Achse erstreckt sich von $-5$ km bis $+5$ km, die y-Achse von $-2.160$ m bis $+540$ m.

\subsection{Dreidimensionale Perspektiven (erosion\_3d.svg)}

Die zweite Visualisierung erweitert die zweidimensionale Darstellung in den dreidimensionalen Raum und ermöglicht eine räumliche Erfassung des Erosionsprozesses. Neun Sub-Plots präsentieren denselben zeitlichen Ablauf wie in der 2D-Darstellung, jedoch aus einer isometrischen Vogelperspektive.

\subsubsection*{Sub-Plot 1: Zyklus 0 (Tag 0) -- Initiale 3D-Ansicht}

Der erste 3D-Querschnitt zeigt eine rechteckige Wasserfläche, die sich über eine Breite von 10 Kilometern (x-Achse: $-5$ km bis $+5$ km) und eine Länge von 20 Kilometern (y-Achse: $0$ km bis $20$ km) erstreckt. Die z-Achse (Höhe) reicht von $-2.160$ m bis $+540$ m.

Die Wasseroberfläche ist als kontinuierliche blaue Fläche (Farbe: \#4A90E2, Transparenz 50\%) dargestellt, die bei $z = 0$ liegt. Unter dieser Oberfläche befindet sich eine braune Gesteinsfläche (Farbe: \#8B7355, Transparenz 70\%), die das Sedimentgestein repräsentiert.

Die Kamera ist unter einem Elevationswinkel von 25° und einem Azimuthwinkel von 45° positioniert, was eine klassische isometrische Ansicht erzeugt. Diese Perspektive ermöglicht es, sowohl die Oberflächenausdehnung als auch die vertikale Dimension gleichzeitig zu erfassen.

Das 3D-Gitter ist in Grau mit 30\% Transparenz dargestellt und ermöglicht eine räumliche Orientierung. Die Achsenbeschriftungen sind in 7-Punkt-Schrift gehalten: „Breite (km)" für die x-Achse, „Länge (km)" für die y-Achse und „Höhe (m)" für die z-Achse.

Der Titel lautet: „Zyklus 0 (Tag 0)", was den Ausgangszustand vor Beginn des Erosionsprozesses bezeichnet. Die Wasseroberfläche erscheint perfekt eben, ohne jegliche Störungen -- ein Zustand des hydraulischen Gleichgewichts.

\subsubsection*{Sub-Plot 2: Zyklus 1 (Tag 5) -- Erste 3D-Erosion}

Der zweite 3D-Querschnitt zeigt die Initiierung der Erosion. Die Wasseroberfläche hat sich auf etwa $z = -200$ m abgesenkt, während sich im Zentrum des Geländes eine beginnende V-förmige Einschnürung abzeichnet.

Die Gesteinsoberfläche ist nicht mehr eben, sondern weist eine Vertiefung auf. Diese Vertiefung erstreckt sich als langgestreckte Rinne in y-Richtung über die gesamten 20 Kilometer Länge. Die maximale Tiefe im Zentrum (bei $x = 0$) beträgt etwa 225 Meter.

Die braune Gesteinsfläche zeigt nun eine dreidimensionale Topographie: Von der Seite ($x = \pm 5$ km) fällt sie kontinuierlich zum Zentrum ab, wo die maximale Erosion stattgefunden hat. Die Querprofile (entlang der x-Achse) zeigen eine parabolische Form, während die Längsprofile (entlang der y-Achse) nahezu flach sind.

Die Wasseroberfläche ist weiterhin als blaue Fläche dargestellt, liegt nun aber etwa 200 Meter unterhalb des ursprünglichen Niveaus. Der Übergang zwischen Wasser und erodiertem Gestein ist deutlich sichtbar.

Die Perspektive (Elevation 25°, Azimuth 45°) ermöglicht einen guten Blick sowohl auf die Erosionsrinne als auch auf die verbleibende Wassermasse. Das 3D-Gitter hilft bei der räumlichen Orientierung und ermöglicht eine Abschätzung der Dimensionen.

\subsubsection*{Sub-Plot 3: Zyklus 2 (Tag 10) -- Zweite 3D-Erosionsphase}

Im dritten 3D-Querschnitt hat sich die Erosionsrinne deutlich vertieft und verbreitert. Die maximale Tiefe im Zentrum beträgt nun etwa 450 Meter, während die Breite der Rinne an der Oberfläche auf etwa 3,6 Kilometer angewachsen ist.

Die Gesteinsoberfläche zeigt eine komplexe dreidimensionale Struktur: Die V-förmige Rinne ist in y-Richtung uniform, was auf einen gleichmäßigen Erosionsprozess entlang der gesamten Länge hindeutet. Die Flanken der Rinne steigen mit einem konstanten Gradienten von etwa 1:8 an.

Die Wasseroberfläche liegt bei etwa $z = -400$ m und bedeckt sowohl die erodierte Rinne als auch die angrenzenden unerodierte Bereiche. Die blaue Färbung mit 50\% Transparenz ermöglicht es, die unterlagernde Topographie zu erkennen.

Die braune Gesteinsfläche zeigt nun zwei erkennbare Terrassierungen -- entsprechend den beiden abgeschlossenen Erosionszyklen. Diese Terrassierungen erscheinen als horizontale Absätze an den Canyon-Flanken und sind besonders gut in der Querrichtung (entlang x) sichtbar.

Die 3D-Perspektive verdeutlicht die enorme räumliche Ausdehnung des Erosionsprozesses: Während die Tiefe auf 450 Meter begrenzt ist, erstreckt sich die Rinne über volle 20 Kilometer in der Länge -- ein Aspekt-Verhältnis von etwa 1:44.

\subsubsection*{Sub-Plot 4: Zyklus 3 (Tag 15) -- Mittlere 3D-Erosion}

Der vierte 3D-Querschnitt dokumentiert die fortschreitende Canyon-Formation. Die Erosionstiefe hat 675 Meter erreicht, was etwa einem Drittel der finalen Tiefe entspricht.

Die dreidimensionale Struktur wird zunehmend komplex: Die Canyon-Flanken weisen nun drei deutlich erkennbare Terrassierungen auf, die den drei abgeschlossenen Erosionszyklen entsprechen. Jede Terrassierung ist etwa 225 Meter über der nächstunteren gelegen und erstreckt sich als horizontale Plattform entlang der gesamten Canyon-Länge.

Die Wasseroberfläche liegt bei etwa $z = -600$ m und zeigt eine leicht wellige Struktur -- ein Artefakt der numerischen Interpolation, das die turbulente Natur der Strömung andeutet. Die blaue Färbung ist transparenter als in den vorherigen Phasen, was die zunehmende Klarheit des ablaufenden Wassers symbolisiert.

Die Gesteinsoberfläche zeigt eine ausgeprägte V-Form im Querschnitt. Die Erosionsbreite an der Oberfläche beträgt etwa 5,4 Kilometer, während sich der Canyon zum Grund auf etwa 0,5 Kilometer verengt. Diese konische Struktur entspricht dem im Hauptdokument beschriebenen Breite-zu-Tiefe-Verhältnis.

Die 3D-Perspektive mit Elevation 25° und Azimuth 45° bietet einen optimalen Blickwinkel, um sowohl die Tiefe als auch die Längenausdehnung des Canyons zu erfassen. Die z-Achse zeigt deutlich die vertikale Stratifizierung, während die x-y-Ebene die horizontale Ausdehnung dokumentiert.

\subsubsection*{Sub-Plot 5: Zyklus 4 (Tag 20) -- Fortgeschrittene 3D-Erosion}

Im fünften 3D-Querschnitt erreicht die Erosion den Halbzeit-Meilenstein von 900 Metern Tiefe. Die Canyon-Struktur ist nun deutlich ausgeprägt und dominiert die Topographie.

Die Gesteinsoberfläche zeigt vier Terrassierungen, die sich als horizontale Bänder an beiden Canyon-Flanken manifestieren. Diese Terrassierungen sind in der 3D-Darstellung besonders eindrucksvoll: Sie erscheinen als Plateaus, die sich über die gesamte Länge des Canyons erstrecken und eine charakteristische Treppenstruktur erzeugen.

Die Wasseroberfläche liegt bei etwa $z = -800$ m und bedeckt den größten Teil des Canyons. Die Wassermasse ist deutlich reduziert im Vergleich zu den frühen Phasen, nimmt aber immer noch ein substantielles Volumen ein.

Die Canyon-Breite an der Oberfläche beträgt etwa 7,2 Kilometer, während die Grundbreite auf etwa 1 Kilometer angewachsen ist. Diese Verbreiterung des Canyon-Grundes ist ein charakteristisches Merkmal des Erosionsprozesses und resultiert aus der zunehmenden lateralen Erosion.

Die 3D-Mesh-Struktur der Gesteinsoberfläche besteht aus 50×50 Gitterpunkten, was eine hinreichend hohe Auflösung für die Darstellung der Topographie bietet. Die Farbcodierung (braun für Gestein, blau für Wasser) ermöglicht eine klare Unterscheidung der verschiedenen Komponenten.

\subsubsection*{Sub-Plot 6: Zyklus 5 (Tag 25) -- Späte 3D-Erosionsphase}

Der sechste 3D-Querschnitt zeigt die späte Erosionsphase mit einer Tiefe von 1.125 Metern. Die Canyon-Struktur ist nun massiv und beeindruckend in ihrer räumlichen Ausdehnung.

Die Gesteinsoberfläche weist fünf deutliche Terrassierungen auf, die eine ausgeprägte Treppenstruktur erzeugen. Jede Terrassierung ist etwa 4 Meter breit (in x-Richtung) und erstreckt sich über die gesamten 20 Kilometer in y-Richtung. Diese Plateaus sind leicht nach außen geneigt, was den hydraulischen Kräften während der Erosion entspricht.

Die Wasseroberfläche liegt bei etwa $z = -1.000$ m und zeigt eine komplexe Topographie mit leichten Wellen und Wirbeln -- ein Hinweis auf die turbulente Strömung während des Rücklaufprozesses. Die blaue Färbung mit 50\% Transparenz ermöglicht es, die unterlagernde Canyon-Struktur durchscheinen zu sehen.

Die Canyon-Breite an der Oberfläche nähert sich dem Maximum von etwa 9 Kilometern. Die Flanken zeigen eine leichte Konvexität, die aus der unterschiedlichen Erosionsresistenz verschiedener Gesteinssorten resultiert -- ein Aspekt, der im Hauptdokument nicht explizit modelliert wurde, aber in der Realität von Bedeutung ist.

Die 3D-Perspektive verdeutlicht die enorme Dimension des Erosionsprozesses: Ein Canyon von 1.125 Metern Tiefe, 9 Kilometern Breite und 20 Kilometern Länge repräsentiert ein erodiertes Volumen von etwa 200 Kubikkilometern -- eine gewaltige Materialmenge.

\subsubsection*{Sub-Plot 7: Zyklus 6 (Tag 30) -- Vorletzte 3D-Phase}

Im siebten 3D-Querschnitt erreicht die Erosion eine Tiefe von 1.350 Metern. Die Canyon-Struktur ist nahezu vollständig ausgebildet und zeigt alle charakteristischen Merkmale der finalen Formation.

Die Gesteinsoberfläche weist sechs Terrassierungen auf, die eine komplexe Treppenstruktur mit ausgeprägten horizontalen und vertikalen Komponenten erzeugen. Die horizontalen Plateaus sind nun deutlich schmaler (etwa 2--3 Meter), während die vertikalen Absätze steiler erscheinen -- ein Resultat der zunehmenden Erosionsintensität in den späteren Zyklen.

Die Wasseroberfläche liegt bei etwa $z = -1.200$ m und bedeckt nur noch den unteren Teil des Canyons. Die blaue Fläche ist deutlich kleiner als in den früheren Phasen und konzentriert sich auf den Canyon-Grund und die unmittelbar angrenzenden Bereiche.

Die Canyon-Breite an der Oberfläche beträgt etwa 10,8 Kilometer, während sich der Canyon zum Grund auf etwa 2 Kilometer verengt. Diese konische Struktur ist in der 3D-Darstellung besonders eindrucksvoll: Die Flanken fallen von den Oberflächen-Rändern steil in die Tiefe ab und erzeugen eine dramatische V-Form.

Die 3D-Perspektive mit Elevation 25° und Azimuth 45° bietet einen optimalen Blickwinkel auf die komplexe Terrassenstruktur. Die Beleuchtung des 3D-Modells (von oben links) erzeugt Schattierungen, die die Tiefe und Struktur des Canyons betonen.

\subsubsection*{Sub-Plot 8: Zyklus 7 (Tag 35) -- Penultimate 3D-Erosion}

Der achte 3D-Querschnitt zeigt den vorletzten Erosionszyklus mit einer Tiefe von 1.575 Metern. Die Canyon-Struktur ist nun nahezu vollständig und entspricht in allen wesentlichen Aspekten der finalen Formation.

Die Gesteinsoberfläche weist sieben Terrassierungen auf, die eine ausgeprägte Stratifizierung erzeugen. Die Terrassierungen sind in der 3D-Darstellung als deutliche horizontale Absätze erkennbar, die sich über die gesamte Canyon-Länge erstrecken. Jede Terrassierung repräsentiert eine Pause zwischen den aktiven Erosionsphasen.

Die Wasseroberfläche liegt bei etwa $z = -1.400$ m und ist als dünne blaue Schicht nur noch im untersten Canyon-Bereich sichtbar. Die Wassermasse ist auf etwa 10\% des ursprünglichen Volumens reduziert, was das bevorstehende Ende des Rücklaufprozesses ankündigt.

Die Canyon-Breite an der Oberfläche erreicht das Maximum von etwa 12,6 Kilometern. Die Canyon-Grundbreite beträgt etwa 3 Kilometer, was einem Breite-zu-Tiefe-Verhältnis am Grund von etwa 1:0,5 entspricht. Diese Geometrie ist charakteristisch für schnelle, katastrophale Erosion.

Die 3D-Mesh-Struktur der Gesteinsoberfläche zeigt eine hohe Detailgenauigkeit mit klar erkennbaren Terrassierungen, Flankenneigungen und Canyon-Grund-Strukturen. Die braune Färbung mit 70\% Transparenz erzeugt einen realistischen Gesteins-Eindruck.

Die räumliche Perspektive verdeutlicht die enormen Dimensionen des Erosionsprozesses: Ein Canyon von fast 1,6 Kilometern Tiefe, über 12 Kilometern Breite und 20 Kilometern Länge repräsentiert ein Volumen, das mit dem realen Grand Canyon vergleichbar ist.

\subsubsection*{Sub-Plot 9: Zyklus 8 (Tag 40) -- 3D-Endzustand}

Der neunte und finale 3D-Querschnitt präsentiert den vollendeten Canyon nach Abschluss aller Erosionszyklen. Die Gesamttiefe beträgt exakt 1.800 Meter -- die dokumentierte Maximaltiefe des Grand Canyon.

Das Wasser ist vollständig abgeflossen; keine blaue Fläche ist mehr vorhanden. Die gesamte Canyon-Struktur liegt offen und bietet einen ungehinderten Blick auf die komplexe dreidimensionale Topographie.

Die Gesteinsoberfläche zeigt acht Hauptterrassierungen, die eine dramatische Treppenstruktur erzeugen. Diese Terrassierungen sind in der 3D-Darstellung besonders eindrucksvoll: Sie manifestieren sich als horizontale Bänder, die sich über die gesamten 20 Kilometer Länge erstrecken und eine charakteristische Schichtung erzeugen, die den realen Grand-Canyon-Formationen bemerkenswert ähnlich ist.

Die Canyon-Breite an der Oberfläche beträgt etwa 14,4 Kilometer -- das absolute Maximum, das durch die akkumulierte Erosion aller Zyklen erreicht wurde. Die Canyon-Grundbreite beträgt etwa 4 Kilometer, was einen substantiellen Flusslauf ermöglichen würde -- konsistent mit der Rolle des Colorado River im modernen Grand Canyon.

Die Flanken des Canyons zeigen eine komplexe Struktur mit ausgeprägten Terrassierungen, steilen vertikalen Absätzen und leicht konkaven Segmenten. Diese Geometrie resultiert aus der Überlagerung von acht diskreten Erosionsphasen, jede mit ihrer charakteristischen Tiefe und lateralen Ausdehnung.

Die 3D-Perspektive (Elevation 25°, Azimuth 45°) bietet einen Panoramablick auf den vollendeten Canyon. Die isometrische Projektion ermöglicht es, sowohl die Tiefe (1,8 km), die Breite (14,4 km) als auch die Länge (20 km) simultan zu erfassen -- ein beeindruckendes Zeugnis der Kraft katastrophaler Erosion.

Das 3D-Gitter mit Transparenz 30\% ermöglicht eine präzise räumliche Orientierung. Die Achsenbeschriftungen geben die Skalen an: x-Achse von $-5$ bis $+5$ km (Breite), y-Achse von $0$ bis $20$ km (Länge), z-Achse von $-2.160$ bis $+540$ m (Höhe).

Der Titel „Zyklus 8 (Tag 40)" kennzeichnet diese Ansicht als finalen Zustand, wobei zu beachten ist, dass „Tag 40" den Beginn der letzten Erosionsphase markiert, die sich bis Tag 220 erstreckt -- eine Phase gradueller Verfestigung und Stabilisierung.

\subsection{Mathematische Formalisierung (erosion\_formulas.svg)}

Die dritte Visualisierung präsentiert das theoretische Fundament des Erosionsmodells in Form mathematischer Gleichungen, die in vier thematische Subsektionen gegliedert sind. Jede Subsektion ist farblich codiert, um die visuelle Unterscheidung zu erleichtern.

\subsubsection*{Oberer linker Quadrant: Hauptgleichungen}

Der erste Quadrant (oben links) trägt den Titel „Hauptgleichungen" in 11-Punkt-Fettschrift und präsentiert die fundamentalen Beziehungen des Erosionsmodells.

Die erste Gleichung definiert die Zyklusperiode:
\begin{equation*}
T_{\mathrm{Zyklus}} = \frac{221 \, \mathrm{Tage}}{45 \, \mathrm{Zyklen}} \approx 4{,}9 \, \mathrm{Tage}
\end{equation*}

Diese Gleichung ist auf hellblauem Hintergrund (Farbe: lightblue, Transparenz 30\%) in monospace-Schrift (10 Punkt) dargestellt. Sie ergibt sich aus der biblischen Chronologie: Die Rücklaufphase dauerte von Tag 150 bis Tag 371, also 221 Tage. Bei angenommenen 45 Erosionszyklen ergibt sich eine mittlere Zyklusdauer von etwa 5 Tagen -- eine bemerkenswerte Übereinstimmung mit der aus den Canyon-Rillenabständen abgeleiteten Periodizität.

Die zweite Gleichung beschreibt die Strömungsgeschwindigkeit:
\begin{equation*}
v = \sqrt{\frac{2 \Delta P}{\rho}} \approx 440 \, \mathrm{m/s}
\end{equation*}

Diese Beziehung folgt aus dem Bernoulli-Prinzip für inkompressible Fluide. Mit einem Druckunterschied $\Delta P \approx 10^7$ Pa (100 bar) und der Wasserdichte $\rho = 1000$ kg/m³ ergibt sich eine Strömungsgeschwindigkeit von etwa 440 m/s -- nahezu Schallgeschwindigkeit und weit über allen natürlichen Flussströmungen.

Die dritte Gleichung quantifiziert den Unterdruck an den Quellen:
\begin{equation*}
\Delta P \approx 10^7 \, \mathrm{Pa} \, (100 \, \mathrm{bar})
\end{equation*}

Dieser extreme Unterdruck ist die treibende Kraft des Rücklaufprozesses. Er entspricht dem hydrostatischen Druck einer Wassersäule von etwa 1.000 Metern Höhe und ist ausreichend, um massive Erosion zu erzeugen.

Die vierte Gleichung berechnet die Schubspannung:
\begin{equation*}
\tau = \frac{1}{2} \rho v^2 \approx 968 \, \mathrm{kPa}
\end{equation*}

Mit $\rho = 1000$ kg/m³ und $v = 440$ m/s ergibt sich eine Schubspannung von etwa 968 kPa -- etwa das 10.000-fache normaler Flussströmungen und mehr als ausreichend, um selbst härteste Gesteine zu erodieren.

Alle Gleichungen sind in separaten Zeilen angeordnet, wobei zwischen den Gleichungen Leerzeilen zur besseren Lesbarkeit eingefügt sind. Jede Gleichung ist in einer abgerundeten Box mit hellblauem Hintergrund präsentiert, was die visuell ansprechende Darstellung unterstützt.

\subsubsection*{Oberer rechter Quadrant: Energiebilanz}

Der zweite Quadrant (oben rechts) trägt den Titel „Energiebilanz" und fokussiert auf die energetischen Aspekte des Erosionsprozesses. Der Hintergrund der Gleichungsboxen ist hellgelb (Farbe: lightyellow, Transparenz 30\%).

Die erste Gleichung definiert die Wassermasse:
\begin{equation*}
m_{\mathrm{Wasser}} = \rho \times V \approx 10^{21} \, \mathrm{kg}
\end{equation*}

Diese Masse ergibt sich aus der Annahme, dass die Sintflut die gesamte Erdoberfläche ($\approx 5 \times 10^{14}$ m²) mit einer mittleren Tiefe von etwa 2.000 Metern bedeckte. Das resultierende Volumen $V \approx 10^{18}$ m³ multipliziert mit der Wasserdichte ergibt die angegebene Masse.

Die nächsten beiden Gleichungen berechnen die potentielle Energie:
\begin{align*}
E_{\mathrm{pot}} &= m \times g \times h \\
E_{\mathrm{pot}} &\approx 10^{21} \times 9{,}8 \times 900 \approx 10^{25} \, \mathrm{J}
\end{align*}

Die mittlere Höhe $h = 900$ m entspricht der Hälfte der angenommenen Wassertiefe. Die potentielle Energie von $10^{25}$ Joule ist eine astronomische Größe -- etwa das 10-fache der jährlichen globalen Energieproduktion der Menschheit und mehr als ausreichend für massive geologische Umgestaltungen.

Die letzten beiden Gleichungen berechnen die kinetische Energie:
\begin{align*}
E_{\mathrm{kin}} &= \frac{1}{2} m v^2 \\
E_{\mathrm{kin}} &\approx \frac{1}{2} \times 10^{21} \times 440^2 \approx 10^{23} \, \mathrm{J}
\end{align*}

Mit der Strömungsgeschwindigkeit $v = 440$ m/s ergibt sich eine kinetische Energie von $10^{23}$ Joule. Diese ist zwar kleiner als die potentielle Energie (etwa 1\% davon), aber immer noch enorm und mehr als ausreichend für die beobachtete Erosion.

Alle Gleichungen sind in 9-Punkt-monospace-Schrift dargestellt und durch Leerzeilen getrennt. Die hellgelbe Hintergrundfarbe unterscheidet diesen Quadranten visuell vom blauen Hauptgleichungs-Quadranten.

\subsubsection*{Unterer linker Quadrant: Erosionsmechanik}

Der dritte Quadrant (unten links) trägt den Titel „Erosionsmechanik" und beschreibt die physikalischen Prozesse der Gesteinserosion. Der Hintergrund der Gleichungsboxen ist hellgrün (Farbe: lightgreen, Transparenz 30\%).

Die erste Gleichung beschreibt die Erosionsrate:
\begin{equation*}
\frac{dh}{dt} = k \times \tau^n
\end{equation*}

Diese empirische Beziehung verbindet die Erosionsrate $dh/dt$ (Tiefenänderung pro Zeit) mit der Schubspannung $\tau$ über eine Konstante $k$ und einen Exponenten $n$ (typischerweise $n \approx 1{,}5$ für Sedimentgestein). Diese Gleichung ist fundamental für quantitative Erosionsmodelle.

Die zweite Gleichung gibt die Schubspannung für Oberflächenströmungen an:
\begin{equation*}
\tau = \rho g h S
\end{equation*}

Hierbei ist $h$ die Wassertiefe und $S$ die Gefälle (slope). Diese Beziehung gilt für Oberflächenströmungen wie Flüsse, ist aber im vorliegenden Modell weniger relevant, da die Erosion primär durch den Unterdruck-Sog getrieben wird, nicht durch Oberflächenströmung.

Die dritte Gleichung definiert den Volumenstrom:
\begin{equation*}
Q = A \times v
\end{equation*}

Der Volumenstrom $Q$ ist das Produkt aus Querschnittsfläche $A$ und Strömungsgeschwindigkeit $v$. Bei einer typischen Canyon-Breite von 10 km und einer Wassertiefe von 100 m ergibt sich $A \approx 10^6$ m². Mit $v = 440$ m/s resultiert ein Volumenstrom von $Q \approx 4 \times 10^8$ m³/s -- ein gewaltiger Wert, der den katastrophalen Charakter des Prozesses unterstreicht.

Die vierte Gleichung berechnet die Reynoldszahl:
\begin{equation*}
Re = \frac{\rho v L}{\mu} \approx 10^{10}
\end{equation*}

Mit $\rho = 1000$ kg/m³, $v = 440$ m/s, $L = 1000$ m (charakteristische Länge) und $\mu = 10^{-3}$ Pa·s (dynamische Viskosität von Wasser) ergibt sich $Re \approx 10^{10}$. Dieser extrem hohe Wert (kritisch ist $Re > 2000$) indiziert hochgradig turbulente Strömung -- ein Regime, in dem die Erosionskraft maximal ist.

Die hellgrüne Hintergrundfarbe kennzeichnet diesen Quadranten als Erosions-Mechanik-Sektion. Die Gleichungen sind in 9-Punkt-monospace-Schrift mit adäquaten Zeilenabständen dargestellt.

\subsubsection*{Unterer rechter Quadrant: Hydraulische Parameter}

Der vierte Quadrant (unten rechts) trägt den Titel „Hydraulische Parameter" und quantifiziert die physikalischen Auswirkungen auf Objekte im Strömungsfeld. Der Hintergrund der Gleichungsboxen ist hellkoralle (Farbe: lightcoral, Transparenz 30\%).

Die ersten beiden Gleichungen berechnen die Sogkraft auf einen Körper:
\begin{align*}
F_{\mathrm{Sog}} &= \Delta P \times A \\
F_{\mathrm{Sog}} &= 10^7 \times 0{,}5 = 5 \times 10^6 \, \mathrm{N}
\end{align*}

Für einen menschlichen Körper mit Querschnittsfläche $A \approx 0{,}5$ m² und Druckunterschied $\Delta P = 10^7$ Pa ergibt sich eine Sogkraft von 5 Millionen Newton -- etwa das Gewicht von 500 Tonnen. Diese Kraft verdeutlicht die absolute Tödlichkeit der Quellen für jedes Lebewesen.

Die nächsten beiden Gleichungen berechnen die Beschleunigung:
\begin{align*}
a &= \frac{F}{m} = \frac{5 \times 10^6}{70} \approx 71{.}000 \, \mathrm{m/s^2}
\end{align*}

Für einen Menschen mit Masse $m = 70$ kg resultiert eine Beschleunigung von etwa 71.000 m/s², was etwa dem 7.200-fachen der Erdbeschleunigung entspricht. Ein zusätzlicher Kommentar in Klammern betont: „(etwa 7.200× Erdbeschleunigung!)".

Die letzte Gleichung berechnet die Machzahl:
\begin{equation*}
M = \frac{v}{c_{\mathrm{Schall}}} = \frac{440}{340} \approx 1{,}3
\end{equation*}

Mit der Strömungsgeschwindigkeit $v = 440$ m/s und der Schallgeschwindigkeit in Luft $c \approx 340$ m/s ergibt sich eine Machzahl von etwa 1,3 -- d.h., die Strömung ist überschallschnell in Bezug auf die Schallgeschwindigkeit in Luft. In Wasser (wo $c \approx 1500$ m/s) ist die Strömung unterschallschnell, aber immer noch extrem schnell.

Die hellkorallfarbene Hintergrundfarbe unterscheidet diesen Quadranten von den anderen drei. Alle Gleichungen sind in 9-Punkt-monospace-Schrift mit konsistenten Zeilenabständen dargestellt.

Die vier Quadranten zusammen bilden eine komprehensive mathematische Beschreibung des Erosionsprozesses, die alle relevanten Aspekte -- von der zeitlichen Periodizität über die Energetik bis zur Erosionsmechanik und den hydraulischen Auswirkungen -- abdeckt.

\subsection{Prozessablauf-Diagramm (erosion\_process.svg)}

Die vierte Visualisierung präsentiert den zeitlichen Ablauf des Erosionsprozesses in Form eines vertikalen Flussdiagramms mit sieben Hauptphasen. Die Darstellung kombiniert grafische Elemente (farbcodierte Rechtecke, Pfeile) mit Textbeschreibungen, um den Prozess in seiner zeitlichen Entwicklung zu dokumentieren.

\subsubsection*{Gesamtstruktur und Layout}

Das Diagramm ist als vertikale Sequenz von sieben Phasen organisiert, die von oben nach unten fortschreitet. Jede Phase wird durch ein horizontal orientiertes Rechteck repräsentiert, das etwa 90\% der Diagrammbreite einnimmt (von x = 0,05 bis x = 0,95) und eine Höhe von etwa 12\% der Gesamthöhe hat.

Die Phasen sind farblich graduiert in verschiedenen Blautönen, die von hellem Himmelblau (oben) zu dunklerem Cyan (unten) übergehen. Die spezifischen Farben sind:
\begin{itemize}
    \item Phase 1: \#E8F4F8 (sehr helles Blau)
    \item Phase 2: \#D0E8F0 (helles Blau)
    \item Phase 3: \#B8DCE8 (mittleres Hellblau)
    \item Phase 4: \#A0D0E0 (mittleres Blau)
    \item Phase 5: \#88C4D8 (mittleres Dunkelblau)
    \item Phase 6: \#70B8D0 (dunkles Blau)
    \item Phase 7: \#58ACC8 (sehr dunkles Blau)
\end{itemize}

Diese Farbgraduierung symbolisiert den fortschreitenden Wasserrückzug: Helle Farben repräsentieren hohe Wasserstände, dunkle Farben niedrige Wasserstände.

Zwischen den Phasen befinden sich vertikale Pfeile (dunkelblau, Kopfbreite 0,03, Kopflänge 0,01), die die zeitliche Progression von einer Phase zur nächsten anzeigen. Die Pfeile sind bei y-Position etwa 0,02 unterhalb jedes Phasen-Rechtecks positioniert.

\subsubsection*{Phase 1: Initialer Wasserstand}

Das erste Rechteck (y-Position von etwa 0,85 bis 0,73) repräsentiert die initiale Phase. Die Beschriftung in der linken Region (x = 0,08) lautet in Fettschrift: „Phase 1". Rechts daneben (x = 0,88, rechtsbündig) steht die Zeitangabe: „Tag 0--30".

Die Prozessbeschreibung in der mittleren Region (x = 0,25, zentriert vertikal) erklärt in kursiver 8-Punkt-Schrift:
\begin{quote}
\textit{„Initialer Wasserstand: Maximum\\
Keine Erosion, Sedimentation abgeschlossen"}
\end{quote}

Diese Phase entspricht dem Zustand am Ende der Sedimentationsphase (biblischer Tag 150). Das Wasser bedeckt die Erde vollständig bis zu einer mittleren Tiefe von etwa 2.000 Metern. Die Sedimente haben sich abgelagert und beginnen zu verfestigen, sind aber noch nicht vollständig konsolidiert. Es findet keine aktive Erosion statt; das System ist in hydraulischem Gleichgewicht.

\subsubsection*{Phase 2: Erster Rücklaufzyklus beginnt}

Das zweite Rechteck (y-Position von etwa 0,73 bis 0,61) markiert den Beginn des katastrophalen Rücklaufs. Die Beschriftung lautet: „Phase 2", Zeitangabe: „Tag 30--60".

Die Prozessbeschreibung erklärt:
\begin{quote}
\textit{„1. Rücklaufzyklus beginnt\\
Sog an Quellen baut sich auf: $\Delta P \approx 10^7$ Pa"}
\end{quote}

In dieser Phase öffnen sich die „Quellen der großen Tiefe" (oder beginnen aktiv Wasser zu absorbieren), was einen extremen Unterdruck von etwa $10^7$ Pa (100 bar) erzeugt. Dieser Unterdruck initiiert den Wasserrückzug und setzt die Strömung in Gang. Die Strömungsgeschwindigkeit beginnt zu steigen und erreicht bald die berechneten 440 m/s.

Das hellblaue Rechteck (Farbe \#D0E8F0) symbolisiert den noch hohen, aber beginnend sinkenden Wasserstand. Die Zeitangabe „Tag 30--60" deutet an, dass diese Phase etwa 30 Tage dauert -- länger als die späteren Zyklen, da das System erst in Gang kommen muss.

\subsubsection*{Phase 3: Massive Erosion}

Das dritte Rechteck (y-Position von etwa 0,61 bis 0,49) repräsentiert die erste aktive Erosionsphase. Beschriftung: „Phase 3", Zeitangabe: „Tag 60--90".

Die Prozessbeschreibung lautet:
\begin{quote}
\textit{„Massive Erosion: $v = 440$ m/s\\
Erste horizontale Rille entsteht"}
\end{quote}

Die Strömungsgeschwindigkeit hat ihren Maximalwert von 440 m/s erreicht, was einer Machzahl von etwa 1,3 (in Bezug auf Luftschall) entspricht. Die damit verbundene Schubspannung von 968 kPa erodiert das frische, noch nicht vollständig verfestigte Sedimentgestein mit extremer Effizienz.

Die erste horizontale Rille entsteht -- eine charakteristische Erosionsstufe, die sich über große Distanzen erstreckt und später als Terrassierung im Canyon sichtbar sein wird. Die Erosionstiefe in dieser Phase beträgt etwa 225 Meter, konsistent mit den in der 2D- und 3D-Visualisierung gezeigten Werten.

Das mittelblaue Rechteck (Farbe \#B8DCE8) symbolisiert den weiter sinkenden Wasserstand. Die 30-Tage-Dauer dieser Phase ist typisch für die frühen bis mittleren Erosionszyklen.

\subsubsection*{Phase 4: Druckausgleich und Pause}

Das vierte Rechteck (y-Position von etwa 0,49 bis 0,37) markiert eine Pausenphase zwischen Erosionszyklen. Beschriftung: „Phase 4", Zeitangabe: „Tag 90--120".

Die Prozessbeschreibung erklärt:
\begin{quote}
\textit{„Druckausgleich, Pause\\
Wasser stabil, keine weitere Erosion"}
\end{quote}

Nach der ersten massiven Erosion erreicht das System einen temporären Gleichgewichtszustand. Der Unterdruck an den Quellen ist vorübergehend ausgeglichen, da die lokale Kapazität der Quellen erschöpft ist. Das globale Quellennetzwerk muss sich neu organisieren und Druckunterschiede ausgleichen, bevor der nächste Erosionszyklus beginnen kann.

Während dieser Pause sinkt der Wasserstand nur minimal oder gar nicht. Die Strömungsgeschwindigkeit reduziert sich auf nahezu null. Das Sediment beginnt zu verfestigen, und die in der vorherigen Phase erodierte Oberfläche stabilisiert sich.

Diese Pausenphase ist entscheidend für die Bildung der horizontalen Rillen: Die stabile Wasseroberfläche über einen Zeitraum von etwa 5 Tagen ermöglicht es, dass sich eine distinkte horizontale Schicht bildet, die später als Terrassierung sichtbar wird.

Das Rechteck in Farbe \#A0D0E0 (mittleres Blau) symbolisiert den stabilen, mittleren Wasserstand. Die Zeitangabe „Tag 90--120" deutet eine 30-Tage-Pause an, was etwa dem Sechsfachen der späteren Zyklusperiode von 5 Tagen entspricht -- ein Hinweis darauf, dass die frühen Phasen längere Ausgleichszeiten benötigten.

\subsubsection*{Phase 5: Zweiter Rücklaufzyklus}

Das fünfte Rechteck (y-Position von etwa 0,37 bis 0,25) beschreibt den zweiten aktiven Erosionszyklus. Beschriftung: „Phase 5", Zeitangabe: „Tag 120--150".

Die Prozessbeschreibung lautet:
\begin{quote}
\textit{„2. Rücklaufzyklus\\
Nächste Erosionsschicht bei tieferem Level"}
\end{quote}

Nach der Pause baut sich der Unterdruck an den Quellen erneut auf, initiiert durch die globale Reorganisation des Quellennetzwerks. Der zweite Erosionszyklus beginnt, nun bei einem niedrigeren Wasserstand als der erste Zyklus.

Die Strömungsgeschwindigkeit erreicht erneut etwa 440 m/s, und eine zweite Erosionsschicht von etwa 225 Metern Tiefe wird erodiert. Diese zweite Schicht liegt etwa 225 Meter unterhalb der ersten, entsprechend dem in Phase 4 erreichten niedrigeren Wasserstand.

Die zweite horizontale Rille entsteht, parallel zur ersten. Der Canyon vertieft sich nun auf etwa 450 Meter Gesamttiefe. Das Erosionsmuster wiederholt sich: intensive Erosion für einige Tage, gefolgt von einer Pause für Druckausgleich.

Das Rechteck in Farbe \#88C4D8 (mittleres Dunkelblau) symbolisiert den weiter gesunkenen Wasserstand. Die Zeitangabe „Tag 120--150" deutet eine 30-Tage-Phase an, konsistent mit den früheren Zyklen.

\subsubsection*{Phase 6: Wiederholende Zyklen}

Das sechste Rechteck (y-Position von etwa 0,25 bis 0,13) fasst die mittleren bis späten Erosionszyklen zusammen. Beschriftung: „Phase 6", Zeitangabe: „Tag 150--350".

Die Prozessbeschreibung erklärt:
\begin{quote}
\textit{„Wiederholende Zyklen\\
Jede Pause = neue horizontale Schicht"}
\end{quote}

Diese Phase repräsentiert die Wiederholung des Erosions-Pausen-Musters über viele Zyklen (etwa 35--45). Jeder Zyklus folgt demselben Schema:
\begin{enumerate}
    \item Aufbau des Unterdrucks an den Quellen
    \item Massive Erosion bei hoher Strömungsgeschwindigkeit ($\approx$ 440 m/s)
    \item Erosion einer etwa 40--50 Meter tiefen Schicht (die Schichten werden dünner in späteren Zyklen)
    \item Druckausgleich und Pause ($\approx$ 5 Tage)
    \item Bildung einer horizontalen Terrassierung
\end{enumerate}

Jede Pause erzeugt eine neue horizontale Schicht, da sich der Wasserstand während der Pause stabilisiert. Die akkumulierte Wirkung vieler Zyklen erzeugt die charakteristische Treppenstruktur mit 40--50 horizontalen Rillen.

Der Wasserstand sinkt kontinuierlich von etwa 1.000 Metern zu Beginn dieser Phase auf etwa 200 Meter am Ende. Die Erosionstiefe wächst von etwa 800 Metern auf etwa 1.600 Meter.

Das Rechteck in Farbe \#70B8D0 (dunkles Blau) symbolisiert den niedrigen Wasserstand. Die Zeitangabe „Tag 150--350" deutet eine lange Phase von 200 Tagen an -- etwa 40 Zyklen à 5 Tage, konsistent mit der im Hauptdokument abgeleiteten Gesamtzahl von 40--50 Zyklen.

\subsubsection*{Phase 7: Finaler Zustand}

Das siebte und letzte Rechteck (y-Position von etwa 0,13 bis 0,01) beschreibt den abschließenden Zustand. Beschriftung: „Phase 7", Zeitangabe: „Tag 350--370".

Die Prozessbeschreibung lautet:
\begin{quote}
\textit{„Finaler Zustand nach $\sim$220 Tagen\\
40--50 Rillen, Tiefe: 1.800 m"}
\end{quote>

Nach etwa 220 Tagen Rücklaufphase (von biblischem Tag 150 bis Tag 371) ist das Wasser vollständig abgeflossen. Der Canyon hat seine finale Tiefe von 1.800 Metern erreicht -- exakt die dokumentierte Maximaltiefe des Grand Canyon.

Die 40--50 Erosionszyklen haben 40--50 horizontale Rillen erzeugt, die sich als charakteristische Terrassierungen über die gesamte Canyon-Länge erstrecken. Diese Rillen sind das definitive geologische Merkmal, das das stufenweise Erosionsmodell vom kontinuierlichen Erosionsmodell unterscheidet.

Die Sedimente sind nun weitgehend verfestigt. Die massive hydraulische Energie ist dissipiert, teilweise in Erosionsarbeit, teilweise in Wärme. Die Landschaft beginnt zu stabilisieren und in die post-Sintflut-Periode überzugehen.

Das Rechteck in Farbe \#58ACC8 (sehr dunkles Blau, fast Cyan) symbolisiert den minimalen Wasserstand oder vollständigen Wasserabfluss. Die Zeitangabe „Tag 350--370" entspricht den letzten 20 Tagen der Sintflut-Chronologie, während derer die Erde „ganz trocken" wurde (Genesis 8:14).

\subsubsection*{Kernmechanismus-Box}

Am unteren Rand des Diagramms (y-Position etwa 0,05, zentriert bei x = 0,5) befindet sich eine hervorgehobene Textbox mit goldenem Hintergrund (Farbe: gold, Transparenz 70\%) und schwarzem Rahmen (Linienbreite 2).

Der Text in dieser Box lautet in 10-Punkt-Fettschrift:
\begin{quote}
\textbf{„Kernmechanismus: Vernetztes Quellensystem mit begrenzter Kapazität\\
→ Rhythmischer Wasserrückzug in Zyklen von $\sim$5 Tagen\\
→ Jede Pause erzeugt eine horizontale Schicht\\
→ 40--50 Zyklen = 40--50 charakteristische Rillen"}
\end{quote}

Diese Box fasst die zentrale Aussage des gesamten Modells zusammen: Die horizontalen Rillen des Grand Canyon sind nicht das Resultat langsamer, kontinuierlicher Erosion, sondern die direkte Manifestation eines im Erdinneren vernetzten, kapazitätsbegrenzten Quellensystems, das einen rhythmischen, stufenweisen Wasserrückzug erzwang.

Die vier Stichpunkte (markiert mit →-Pfeilen) erklären die kausale Kette:
\begin{enumerate}
    \item Das globale Quellennetzwerk hatte begrenzte Kapazität
    \item Dies erzwang rhythmischen Rückzug in Zyklen von etwa 5 Tagen
    \item Jede Pause zwischen Zyklen stabilisierte das Wasser und erzeugte eine horizontale Schicht
    \item Die Summe von 40--50 solcher Zyklen erzeugte 40--50 charakteristische Rillen
\end{enumerate}

Diese Box ist visuell hervorgehoben durch den goldenen Hintergrund und die zentrale Position am Ende des Diagramms. Sie dient als abschließende Zusammenfassung und „Take-Home-Message" der gesamten Visualisierung.

Die vier Visualisierungen zusammen -- 2D-Querschnitte, 3D-Perspektiven, mathematische Formeln und Prozessablauf -- bilden eine komprehensive Darstellung des theoretischen Wasserabsaug- und Erosionsprozesses, der zur Entstehung des Grand Canyon führte. Sie kombinieren räumliche, zeitliche und mathematische Perspektiven, um ein vollständiges Bild des postulierten katastrophalen Prozesses zu vermitteln.

